%\begin{previewactivity}[The Telephone Keypad] \label{PA:telephone} \hfill
%
%The following diagram represents the keypad for a telephone.
%
%\begin{figure}[h]
%\begin{center}
%\includegraphics{figps-phone.eps}
%%\caption{Composition of Functions} \label{fig:functioncomposition2}
%\end{center}
%\end{figure}
%
%Let  $\Gamma  = \left\{ {0, 1, 2,  \ldots , 8, 9} \right\}$, and let  $\Delta $ be the set of all upper case letters in the English alphabet.  (\underline{Note}:  We are using the uppercase Greek letters gamma, $\Gamma $, and  delta, $\Delta $, to avoid confusion with the letters on the keypad.)  Define  $\Omega $, the uppercase Greek letter omega, to be the following subset of  
%$\Gamma  \times \Delta $:
%\[
%\Omega  = \left\{ { {\left( {x, y} \right) \in \Gamma  \times \Delta  } \mid x\text{  and  }y\text{  are on the same button}} \right\}.
%\]
%
%\begin{enumerate}
%\item Determine at least five different ordered pairs in  $\Gamma  \times \Delta $  that are elements of  $\Omega $.
%
%\item Determine at least five different ordered pairs in  $\Gamma  \times \Delta $  that are not elements of  $\Omega $.
%
%\item Determine all elements of the set  
%$\left\{ {\ {x \in \Gamma  } \mid \left( {\exists y \in \Delta } \right)\left( {\left( {x, y} \right) \in \Omega } \right)} \right\}$. \label{PA:telephone3}
%
%\item Determine all elements of the set  
%$\left\{ { {y \in \Delta  } \mid \left( {\exists x \in \Gamma } \right)\left( {\left( {x, y} \right) \in \Omega } \right)} \right\}$.  \label{PA:telephone4}
%\end{enumerate}
%\end{previewactivity}
%\hbreak
%

\begin{previewactivity}[The United States of America] \label{PA:USA} \hfill \\
Recall from Section~\ref{S:cartesian} that the \textbf{Cartesian product}
\index{Cartesian product}%
 of two sets  $A$  and  $B$, written  $A \times B$, is the set of all ordered pairs  
$\left( {a,b} \right)$,  where  $a \in A$  and  $b \in B$.  That is,
$A \times B = \left\{ {\left( {a,b} \right)  \mid a \in A\text{ and }b \in B} \right\}$.

\vskip6pt
Let  $A$  be the set of all states in the United States and let
\[
R = \left\{ { {\left( {x, y} \right) \in A \times A } \mid x \text{  and  }y 
\text{  have a land border in common}} \right\}\!.
\]
\begin{enumerate}
\item Use the roster method for defining a set to list all the elements in each of the following sets:
\begin{enumerate}
\item $A = \left\{ {y \in A\left| {\left( {\text{Michigan, }y} \right) \in R} \right.} \right\}$

\item $B = \left\{ {x \in A\left| {\left( {x,\text{Michigan}} \right) \in R} \right.} \right\}$

\item $C = \left\{ {y \in A\left| {\left( {\text{Wisconsin, }y} \right) \in R} \right.} \right\}$

\end{enumerate}

\item Find two different examples of two ordered pairs,  $\left( {x, y} \right)$ and 
$\left( {y, z} \right)$ such that  $\left( {x, y} \right) \in R$,  
$\left( {y, z} \right) \in R$,  but  $\left( {x, z} \right)\not  \in R$, or explain why no such example exists.

%\item Are the following statements true or false?  Justify your conclusions.
%\begin{enumerate}
%\item For all $x, y \in A$, if $(x, y) \in R$, then $(y, x) \in R$.
%\item For all $x, y, z \in A$, if $(x, y) \in R$ and $(y, z) \in R$, then $(x, z) \in R$.
%\end{enumerate}

\item In Section~\ref{S:inversefunctions}, we learned how to represent a function as a set of ordered pairs, and we learned under what conditions a set of ordered pairs can be used to define a function. Can the set  $R$  be used to define a function from the set  $A$  to the set  $A$?  Explain.
\end{enumerate}
\end{previewactivity}
\hbreak

\begin{previewactivity}[An Equation with Two Variables] \label{PA:eqn2variables} \hfill \\
In Section~\ref{S:predicates}, we introduced the concept of the \textbf{truth set of a predicate with one variable}.
\index{truth set}%
  This was defined to be the set of all elements in the universal set that can be substituted for the variable to make the predicate a true proposition.

Assume that  $x$  and  $y$  represent real numbers.  Then 

\[
4x^2  + y^2  = 16
\]
is a predicate with two variables.  An element of the truth set of this predicate (also called a solution of the equation) is an ordered pair  $\left( {a, b} \right)$ of real numbers so that when  $a$  is substituted for  $x$  and  $b$  is substituted for  $y$, the predicate becomes a true proposition (a true equation in this case).  We can use set builder notation to describe the truth set $S$ of this equation with two variables as follows:
\[
S = \left\{ (x, y) \in \R \times \R \mid 4x^2 + y^2 = 16 \right\}\!.
\]
When a set is a truth set of a predicate that is an equation, we also call the set the 
\textbf{solution set} of the equation.
\begin{enumerate}
\item List four different elements of the set $S$\!.

\item The graph of the equation  $4x^2  + y^2  = 16$  in the $xy$-coordinate plane is an ellipse.  Draw the graph and explain why this graph is a representation of the truth set (solutions set) of the equation $4x^2 + y^2 = 16$.

%\item Write a description of the solution set  $S$  of the equation  $x^2  + y^2  = 25$ using set builder notation.

%\pagebreak
\item Describe each of the following sets as an interval of real numbers:
\label{PA:eqn2variables3}
\begin{enumerate}
\item $A = \left\{ x \in \R \mid \text{ there exists a } y \in \R \text{ such that } 4x^2 + y^2 = 16 \right\}$.

\item $B = \left\{ y \in \R \mid \text{ there exists an } x \in \R \text{ such that } 4x^2 + y^2 = 16 \right\}$.
\end{enumerate}

\end{enumerate}
\end{previewactivity}
\hbreak
%

\begin{previewactivity}[A Set of Order Pairs] \label{PA:orderedpairs} \hfill \\
Let $\R^* = \left\{ y \in \R \mid y \geq 0 \right\}$ and let
$F = \left\{ (x, y) \in \R \times \R^* \mid y = x^2 \right\}$.

\begin{enumerate}
\item List five different ordered pairs that are in the set $F$.

\item Use the roster method to specify the elements of each of the following the sets:
\begin{multicols}{2}
\begin{enumerate}
\item $A = \left\{ x \in \R \mid (x, 4) \in F \right\}$
\item $B = \left\{ x \in \R \mid (x, 10) \in F \right\}$
\item $C = \left\{ y \in \R^* \mid (5, y) \in F \right\}$
\item $D = \left\{ y \in \R^* \mid (-3, y) \in F \right\}$
\end{enumerate}
\end{multicols}

\item Can the set  $F$  be used to define a function from the set  $\R$  to the set  $\R^*$?  Explain.

\end{enumerate}
\end{previewactivity}
\hbreak
\endinput


\begin{center}
\setlength{\unitlength}{0.5cm}
\begin{picture}(16,13)
\put(1,10){\line(1,0){4}}
\put(1,10){\line(0,1){2}}
\put(1,12){\line(1,0){4}}
\put(5,10){\line(0,1){2}}
\put(2.8,11){1}

\put(6,10){\line(1,0){4}}
\put(6,10){\line(0,1){2}}
\put(6,12){\line(1,0){4}}
\put(10,10){\line(0,1){2}}
\put(6.8,11.3){A B C}
\put(7.8,10.2){2}

\put(11,10){\line(1,0){4}}
\put(11,10){\line(0,1){2}}
\put(11,12){\line(1,0){4}}
\put(15,10){\line(0,1){2}}
\put(11.8,11.3){D E F}
\put(12.8,10.2){3}

\put(1,7){\line(1,0){4}}
\put(1,7){\line(0,1){2}}
\put(1,9){\line(1,0){4}}
\put(5,7){\line(0,1){2}}
\put(2,8.3){G H I}
\put(2.8,7.2){4}

\put(6,7){\line(1,0){4}}
\put(6,7){\line(0,1){2}}
\put(6,9){\line(1,0){4}}
\put(10,7){\line(0,1){2}}
\put(6.9,8.3){J K L}
\put(7.8,7.2){5}

\put(11,7){\line(1,0){4}}
\put(11,7){\line(0,1){2}}
\put(11,9){\line(1,0){4}}
\put(15,7){\line(0,1){2}}
\put(11.8,8.3){M N O}
\put(12.8,7.2){6}

\put(1,4){\line(1,0){4}}
\put(1,4){\line(0,1){2}}
\put(1,6){\line(1,0){4}}
\put(5,4){\line(0,1){2}}
\put(1.6,5.3){P Q R S}
\put(2.8,4.2){7}

\put(6,4){\line(1,0){4}}
\put(6,4){\line(0,1){2}}
\put(6,6){\line(1,0){4}}
\put(10,4){\line(0,1){2}}
\put(6.9,5.3){T U V}
\put(7.8,4.2){8}

\put(11,4){\line(1,0){4}}
\put(11,4){\line(0,1){2}}
\put(11,6){\line(1,0){4}}
\put(15,4){\line(0,1){2}}
\put(11.5,5.3){W X Y Z}
\put(12.8,4.2){9}


\put(1,1){\line(1,0){4}}
\put(1,1){\line(0,1){2}}
\put(1,3){\line(1,0){4}}
\put(5,1){\line(0,1){2}}
\put(2.8,2){*}

\put(6,1){\line(1,0){4}}
\put(6,1){\line(0,1){2}}
\put(6,3){\line(1,0){4}}
\put(10,1){\line(0,1){2}}
\put(7.8,2){0}

\put(11,1){\line(1,0){4}}
\put(11,1){\line(0,1){2}}
\put(11,3){\line(1,0){4}}
\put(15,1){\line(0,1){2}}
\put(12.8,2){\#}

\end{picture}
\end{center}
